\documentclass[a4paper,11pt]{article}
\usepackage[margin=27mm]{geometry}
\usepackage{fontspec,amsmath,amssymb,amsthm,unicode-math}
\setmainfont{STIXGeneral}\setmathfont{STIX Two Math}
\setlength{\emergencystretch}{2em}
\newcommand{\Sym}{\operatorname{Sym}}\newcommand{\Alt}{\operatorname{Alt}}
\newcommand{\gen}[1]{\langle#1\rangle}
\newtheorem{theorem}{Theorem}\newtheorem{lemma}[theorem]{Lemma}
\newtheorem{corollary}[theorem]{Corollary}
\title{Restriction kernels in extensions of string markings}
\author{Human Compression project}\date{Private theorem development; independent review pending}
\begin{document}\maketitle
\begin{abstract}
We characterize when adjoining a transposition at a new point preserves a
string C-marking. The tests are inclusions between pointwise stabilizers in
successive original prefix groups. They do not assume the intersection
property of the enlarged marking. Applying the criterion on the two sides
of a split gives a necessary and sufficient condition for the rank and
degree extension when one internal split factor is trivial. A parity
example explains the obstruction. The last end-intersection in the printed
proof of Cameron--Fernandes--Leemans can be established directly; the missing
condition is the intersection property of its recursively obtained ends.
No conclusion about the source's main stability theorem is inferred.
\end{abstract}

\section{Conventions and the source of the obstruction}
All permutation sets are finite. The notation $\Sym(X)$ denotes permutations
supported on $X$, with every point outside $X$ fixed. For a group preserving
$X$, its restriction to the complement denotes the actual induced image,
extended by the identity on $X$. A string C-marking consists of distinct
nonidentity involutions, with distant marks commuting and with
\[
 \gen{s_j:j\in I}\cap\gen{s_j:j\in J}=\gen{s_j:j\in I\cap J}
\]
for all label subsets. Every subfamily of a C-marking has this property.
We retain the standard end criterion: a string marking is C if both end
deletions are C and their groups intersect in the middle group. This is
McMullen--Schulte, Proposition 2E16, recalled as Proposition 2.1 of [CFL].

The exact source [CFL] is arXiv:2212.12723v2, 11 January 2023. Its
Proposition 2.4 asserts an invalid unrestricted converse for removing an
endpoint central twist. Its Theorem 6.4 also fails as stated, even for a
transitive input. The separate complete counterexamples are not repeated
here. The proposed repair concerns the additional group information needed
by the construction, not a change in the meaning of its source hypotheses.
The pointwise stabilizers below detect correlations which disappear when
an induced permutation action is mistaken for an independently supported
subgroup.

\section{Adjoining one transposition}
\begin{lemma}\label{lem:factor}
Let $P\leq\Sym(\Omega)$, let $a\in\Omega$, and let $c$ be a new point.
Extend $P$ to fix $c$, put $X=a^P$ and $t=(a\ c)$. Then
\[
 \gen{P,t}=\Sym(X\cup\{c\})\times P^{\Omega\setminus X}.
\]
\end{lemma}
\begin{proof}
Conjugating $t$ by $P$ gives every $(x\ c)$ for $x\in X$. These
transpositions generate the full symmetric group on $X\cup\{c\}$:
$(x\ c)(y\ c)(x\ c)=(x\ y)$. Cancel the restriction to $X$ of
any element of $P$ using that supported symmetric group. What remains
is its actual restriction to $\Omega\setminus X$, fixing the other
points. Thus both displayed factors are contained in the generated
group. Conversely the original generators preserve their supports and
have the indicated complementary restrictions. The supports are disjoint,
so the product is direct.
\end{proof}

\begin{lemma}\label{lem:kernel}
Let $P\leq Q\leq\Sym(\Omega)$ and $X=a^P$. Then
\[
 Q\cap\gen{P,(a\ c)}=P
 \quad\Longleftrightarrow\quad Q\cap\Sym(X)\leq P.
\]
The group $Q\cap\Sym(X)$ is the pointwise stabilizer in $Q$ of
$\Omega\setminus X$.
\end{lemma}
\begin{proof}
An element of $Q$ fixes $c$. By Lemma~\ref{lem:factor}, its membership
in the second group on the left means that it preserves $X$ and that
its complementary restriction belongs to $P^{\Omega\setminus X}$.
Choose $p\in P$ with that same restriction. Multiplying by $p^{-1}$
leaves an element of $Q$ supported on $X$. The right-hand inclusion
therefore implies that the original element belongs to $P$. The reverse
containment is automatic. Conversely every element of $Q\cap\Sym(X)$
belongs to $Q\cap\gen{P,(a\ c)}$, so the equality implies the inclusion.
\end{proof}

\begin{theorem}\label{thm:leaf}
Let $(h_0,\ldots,h_{k-1})$ be a string C-marking on $\Omega$, and
suppose $h_j$ fixes $a$ for every $j>0$. Put
\[
 P_j=\gen{h_0,\ldots,h_{j-1}},\qquad X_j=a^{P_j}
 \quad(0\leq j\leq k),
\]
where $P_0=1$. For a new point $c$, the marking
\[
 ((a\ c),h_0,\ldots,h_{k-1})
\]
is C if and only if
\begin{equation}\label{eq:kerneltest}
 P_j\cap\Sym(X_{j-1})\leq P_{j-1}
 \qquad(1\leq j\leq k).
\end{equation}
The first condition is automatic. For $k=0$ the enlarged marking has rank one.
\end{theorem}
\begin{proof}
The new transposition commutes with $h_j$ for $j>0$, since each of
those involutions fixes both its points. It is distinct from every old
mark because it moves the new point.

Assume \eqref{eq:kerneltest}. Induct on $j$ to show that
$((a\ c),h_0,\ldots,h_{j-1})$ is C. The initial rank-one case is
immediate. At the next stage one end is the old prefix with group $P_j$,
which is C by hypothesis. The other is the preceding extended prefix,
which is C by induction. Their intersection is $P_{j-1}$ by
Lemma~\ref{lem:kernel} and the $j$th condition. The end criterion
proves the assertion.

Conversely, if the entire new marking is C, every indicated prefix is C.
Its end groups must intersect in $P_{j-1}$. Lemma~\ref{lem:kernel}
gives \eqref{eq:kerneltest} for each $j$. Since $X_0=\{a\}$ and
$\Sym(\{a\})=1$, the first condition is indeed automatic.
\end{proof}

For $k\geq1$, these are at most $k-1$ nontrivial tests in groups generated by the original
marks. Each is a pointwise-stabilizer computation followed by containment;
none asks for an intersection property in the enlarged degree. Equivalently,
the two pointwise stabilizers
\[
 P_j{}_{(\Omega\setminus X_{j-1})}
 \quad\hbox{and}\quad
 P_{j-1}{}_{(\Omega\setminus X_{j-1})}
\]
must be equal. For finite groups equality can be checked by their orders,
because the second is a subgroup of the first.

\begin{corollary}\label{cor:pure}
Under the hypotheses of Theorem~\ref{thm:leaf}, if $h_0=(a\ b)$
is a pure transposition, adjoining $(a\ c)$ preserves the C-property.
\end{corollary}
\begin{proof}
For each $j\geq1$, the subgroup generated by $h_1,\ldots,h_{j-1}$
fixes $a$. Apply Lemma~\ref{lem:factor} with $a$ as its new point
and $b$ as the distinguished old point. It follows that
$P_j$ contains $\Sym(X_j)$. Hence every group on the left of
\eqref{eq:kerneltest} is contained in $\Sym(X_{j-1})\leq P_{j-1}$,
with the trivial $j=1$ case treated separately.
\end{proof}

\section{Two independent sides and a central twist}
For disjoint sets $A,B,\{c\}$, points $a\in A,b\in B$ and groups
$L\leq\Sym(A)$, $T\leq\Sym(B)$, write
\[
 K(L,T)=\gen{L,T,(a\ c),(c\ b)}.
\]
\begin{lemma}\label{lem:rectangle}
If $L_0\leq L_1$ and $T_0\leq T_1$, then
\[
 K(L_0,T_1)\cap K(L_1,T_0)=K(L_0,T_0).
\]
\end{lemma}
\begin{proof}
Writing $X=a^L$ and $Y=b^T$, conjugation of the two pure transpositions
and cancellation of restrictions, as in Lemma~\ref{lem:factor}, give
\[
 K(L,T)=\Sym(X\cup Y\cup\{c\})
 \times L^{A\setminus X}\times T^{B\setminus Y}.
\]
Indeed, the star transpositions from both sides generate the first factor;
the remaining restrictions are separately available on disjoint sets.
Put $X_j=a^{L_j}$ and $Y_j=b^{T_j}$. The two enlarged root sets
intersect in $Z=X_0\cup Y_0\cup\{c\}$. An element of the intersection
preserves $Z$, acts on $A\setminus X_0$ in the actual restriction image
of $L_0$, and acts on $B\setminus Y_0$ in that of $T_0$. These three
sets partition the whole domain, so the displayed formula puts it in
$K(L_0,T_0)$. The reverse inclusion follows from the two containments.
\end{proof}

We also need the valid forward central-extension fact, and include its
proof to avoid invoking the invalid converse. This forward statement is
Fernandes--Leemans--Mixer (2012), Proposition 5.3, as cited in [CFL,
Proposition 2.2] and the gluing source.
\begin{lemma}\label{lem:twist}
Let $(s_0,\ldots,s_{m-1})$ be a string C-marking of $H$. In an ambient
direct product $H\times\gen{z}$ with $z$ of order two, replace just the
last mark by $s_{m-1}z$. The resulting marking is C. The same holds for
$z=1$ and for the first endpoint.
\end{lemma}
\begin{proof}
Projection to $H$ sends each new mark to its old counterpart, preserving
distinctness and nonidentity. The string commutations are unchanged.
Induct on $m$. The rank-one case is immediate. One end deletion is
untwisted and C; the other is an endpoint twist of a shorter C-marking
and is C by induction. An element in their intersection projects into
the old middle subgroup by the old C-property. Projection is the identity
on the untwisted end, so that element itself is its projected middle
element. All middle marks are untwisted. The two new ends therefore
intersect precisely in their middle group, and the end criterion finishes
the induction. Reversal treats the first endpoint.
\end{proof}

\section{A criterion for the split extension}
Let $(\rho_0,\ldots,\rho_{r-1})$ be a string C-marking on
$\Omega=A\sqcup B$, with $a\in A,b\in B$, such that
\begin{equation}\label{eq:split}
 \rho_j\text{ is supported on }A\ (j<i),\qquad
 \rho_j\text{ is supported on }B\ (j>i),\qquad
 \rho_i=(a\ b)\beta,
\end{equation}
where $\beta$ is an involution or the identity supported on $B\setminus\{b\}$.
This is the perfect-split form with $\alpha_i=1$; we do not require
transitivity, a fracture graph, or moved endpoints in the following criterion.
Introduce a point $c$ and replace $\rho_i$ by
\[
 u=(a\ c),\qquad v=(c\ b)\beta.
\]
Consider the following two original string submarkings, both rooted at $a$:
\[
 \mathcal L=(\rho_{i-1},\ldots,\rho_0),\qquad
 \mathcal R=(\rho_i,\rho_{i+1},\ldots,\rho_{r-1}).
\]
The first may be empty. Apply the prefix notation of
Theorem~\ref{thm:leaf} separately to each.

\begin{theorem}\label{thm:split}
The prescribed extended marking is C if and only if both
$\mathcal L$ and $\mathcal R$ satisfy \eqref{eq:kerneltest} at $a$.
There are at most $r$ tests on prefix kernels, at most $r-2$ of them
nontrivial when both lists are nonempty. The symmetric construction in
which the internal factor on $B$ is trivial is obtained by reversing
the order and interchanging the two sides.
\end{theorem}
\begin{proof}
All marks of $\mathcal R$ after its first fix $a$. For $j<i-1$,
commutation of $\rho_j$ with $\rho_i$ forces $\rho_j$ to fix $a$:
evaluate the equality at $b$, which every left mark fixes. Thus all
marks of $\mathcal L$ after its first also fix $a$.

The inserted marks are nonidentity involutions, distinct from each other
and from the old marks because they move the new point.
The extended marking satisfies the string commutations. To verify the
only potentially new ones, commutation of an old right mark with
$\rho_i$ at distance at least two forces it to fix $b$ and to commute
with $\beta$. The old left marks at distance at least two fix $a$,
and they all commute with $\beta$ by disjoint support. These statements
give every distant pair involving $u$ or $v$.

The kernel tests and Theorem~\ref{thm:leaf} make the left marking
$(\rho_0,\ldots,\rho_{i-1},u)$ C, after reversal. They also make
$(u,\rho_i,\rho_{i+1},\ldots)$ C. Conjugating the latter by $u$
fixes $u$ and every old right mark, and sends $\rho_i$ to $v$.
Thus the entire right marking $(u,v,\rho_{i+1},\ldots)$ is C.

We verify every consecutive interval of the extended marking by induction
on its length. Intervals not containing both $u,v$ are submarkings of
the left or right C-markings just identified, or old submarkings, so
are C. The same is true of intervals starting at $u$ and extending
rightwards. An interval ending at $v$ is C as well: first extend its
left part, ending in the pure mark $u$, by the pure $(c\ b)$, using
Corollary~\ref{cor:pure} after reversal. Then multiply just this last
mark by $\beta$. Its support is disjoint from the entire preceding
construction, so Lemma~\ref{lem:twist} applies.

It remains to treat an interval having at least one old mark on each
side of $u,v$. Its end deletions are C by induction. Moreover
$(uv)^3=\beta$, since its three-cycle and the disjoint $\beta$
commute. Hence a selected subgroup containing both $u,v$ is exactly
$K(L,T)$, where $L$ is its selected left subgroup and $T$ is generated
by $\beta$ and its selected old right marks. Deleting the first or
last mark changes these to $K(L_0,T_1)$ and $K(L_1,T_0)$, with
$L_0\leq L_1$, $T_0\leq T_1$. Lemma~\ref{lem:rectangle} identifies
their intersection with the middle group. The end criterion completes
the interval induction and proves the whole marking C.

Conversely, a C extended marking has both displayed left and right
submarkings C. Reverse the first and conjugate the second by $u$.
The necessity direction of Theorem~\ref{thm:leaf} then gives exactly
the two families of tests on original prefix kernels.
\end{proof}

\begin{corollary}\label{cor:subdivide}
If $\rho_i=(a\ b)$ is itself pure in \eqref{eq:split}, replacing it
by $(a\ c),(c\ b)$ preserves the string C-marking.
\end{corollary}
\begin{proof}
The old left submarking ending in $\rho_i$ is C; renaming its point
$b$ as $c$ gives the left marking needed in the proof above. The
right submarking begins with the pure $\rho_i$ and every later mark
fixes $a$, so Corollary~\ref{cor:pure} gives the right marking. The
proof of Theorem~\ref{thm:split} therefore applies.
\end{proof}

\section{Parity fibers and explicit uses}
The criterion can be evaluated from a subgroup description without a
full subgroup census. Suppose $P$ preserves $X$ and $Y=\Omega\setminus X$
and has the fiber form
\[
 P=\{(g,d)\in L\times D:\phi(g)=\psi(d)\},
\]
where $L\leq\Sym(X)$, $D\leq\Sym(Y)$ and $\phi,\psi$ are
homomorphisms to a common group. Then
$P\cap\Sym(X)=\ker\phi$, viewed as a supported subgroup.
For $P\leq Q$, Lemma~\ref{lem:kernel} asks exactly that
$Q\cap\Sym(X)$ lie in this kernel. The quotient maps in such a
description are the familiar fiber-product data of Goursat's lemma;
no classification by that lemma is needed for this direct computation.

In the transitive counterexample the relevant original suffix is
\[
 q=(12)(34),\quad r=(26)(35)(47),\quad s=(34)(67).
\]
It is a C-marking of $S_7$, as proved in the separate note. For the
prefix $P_2=\gen{q,r}$, its root orbit at $1$ is $X_2=\{1,2,6\}$.
Put $t=(34)$, $w=(35)(47)$ and $D=\gen{t,w}\cong D_8$.
The written subgroup formula there is
\[
 P_2=\{g d:g\in\Sym(X_2),\ d\in D,\ \operatorname{sgn}(g)=\chi(d)\},
\]
where $\chi(t)=\chi(w)$ is the nontrivial element of $C_2$.
Consequently $P_2\cap\Sym(X_2)=\Alt(X_2)$, whereas
$P_3\cap\Sym(X_2)=\Sym(X_2)$ because $P_3=S_7$.
The third prefix test fails with index two. This is the precise lost
parity relation, detected in original groups of degrees at most seven.

The test is not a requirement that every original subgroup be an independent
direct product. For example, on points $v_0,\ldots,v_k$ and a disjoint
pair interchanged by $z$, take
\[
 h_j=(v_j\ v_{j+1})z\qquad(0\leq j<k).
\]
This is the image of the adjacent-transposition marking of $S_{k+1}$
under the injective homomorphism $g\mapsto g z^{\varepsilon(g)}$,
so is a string C-marking. Every prefix has this same parity-graph form.
At step $j\geq2$, its kernel on the complement of
$X_{j-1}=\{v_0,\ldots,v_{j-1}\}$ is precisely $\Alt(X_{j-1})$,
equal to the preceding prefix's kernel. Thus a pure new leaf at $v_0$
preserves C even though the original first mark is not pure. Subsequent
pure leaves may be added repeatedly by Corollary~\ref{cor:pure}.
These are written uniform deductions, not extrapolations from sample orders.

For the already admitted gluing families, any interior pure edge of the
added tail has precisely the separated lower and higher supports of
\eqref{eq:split}. Corollary~\ref{cor:subdivide} therefore justifies its
subdivision uniformly, with the folded or doubled core left intact.
This is a bounded application of the repaired theorem, not a new
classification or a proof of all the high-rank cases in [CFL].

\section{What is and is not repaired}
When both inserted marks occur, their cube already separates $\beta$.
The product calculation and the end-intersection on p.~33 of [CFL]
can therefore be supplied by Lemma~\ref{lem:rectangle}, without an
unproved independence of restriction images. That equality alone was
not sufficient: the end criterion also needs both ends to be C, and
recursion reaches endpoint attachments at impure marks. Those attachments
are governed by Theorem~\ref{thm:leaf}; the old C-property and a
fracture graph do not force its kernel inclusions. This is why adding
transitivity to the printed hypotheses does not repair Theorem 6.4.

The standard end criterion and forward twist theorem are attributed above.
The supported-transposition argument reuses the common gluing proof; the
new interface isolates the exact additional restriction-kernel condition
and its two-sided use. No historical priority assertion is made.
Verification of these conditions throughout the main stability range,
and recovery of the remaining classification and counting arguments of
[CFL], remain separate obligations. The main stability theorem is not
claimed false.

The accompanying Lean file verifies three general subgroup statements:
the cancellation argument for equal images under a homomorphism, its
equivalence with the intersection of a subgroup and an inverse image,
and the equivalent equality of kernels. These are actual subgroup
statements, valid without finiteness. The permutation restriction maps,
the supported symmetric factors and the induction using the end criterion
are proved above; they are not formalized by those three declarations.
The complete GAP file tests the criterion and its hypotheses in bounded
examples. Its finite outputs are diagnostics, not premises of the
quantified theorems.
\begin{thebibliography}{9}
\bibitem{CFL} P. J. Cameron, M. E. Fernandes and D. Leemans,
\emph{The number of string C-groups of high rank},
arXiv:2212.12723v2, 11 January 2023. Exact bound source edition;
Propositions 2.1--2.4 and Theorem 6.4 are the locators discussed here.
\end{thebibliography}
\end{document}
